\subsection{ACTIONS}

DEFINITION 1. Let \(\Omega\) and \(\mathbf{E}\) be two sets. A mapping of \(\Omega\) into the set \(\mathbf{E}^{\mathbf{E}}\) of mappings of \(\mathbf{E}\) into itself is called an action of \(\Omega\) on \(\mathbf{E}\) .

Let \(\alpha \mapsto f_{\alpha}\) be an action of \(\Omega\) on E. The mapping \((\alpha, x) \mapsto f_{\alpha}(x)\) (resp. \((x, \alpha) \mapsto f_{\alpha}(x)\) ) is called the law of left (resp. right) action of \(\Omega\) on E† associated with the given action of \(\Omega\) on E. Given a mapping \(g\) of \(\Omega \times E\) (resp. E × \(\Omega\) ) into E, there exists one and only one action \(\alpha \mapsto f_{\alpha}\) of \(\Omega\) on E such that the associated law of left (resp. right) action is \(g\) (Set Theory, II, § 5, no. 2, Proposition 3).

In this chapter we shall say, for the sake of abbreviation, ``law of action''

instead of ``law of left action''. The element \(f_{\alpha}(x)\) of E (for \(\alpha \in \Omega\) and \(x \in E\) ) is sometimes called the transform of x under \(\alpha\) or the composition of \(\alpha\) and x. It is often denoted by left (resp. right) multiplicative notation \(\alpha.x\) (resp. \(x.\alpha\) ), the dot may be omitted; the composition of \(\alpha\) and x is then called the product of \(\alpha\) and x (resp. x and \(\alpha\) ). The exponential notation \(x^{\alpha}\) is also used. In the arguments of the following paragraphs we shall generally use the notation \(\alpha \perp x\) . The elements of \(\Omega\) are often called operators.

Examples. (1) Let E be an associative magma written multiplicatively. The mapping which associates with a strictly positive integer n the mapping \(x \mapsto x^{n}\) of E into itself is an action of \(N^{*}\) on E. If E is a group, the mapping which associates with a rational integer a the mapping \(x \mapsto x^{a}\) of E into E is an action of Z on E.

\begin{enumerate}
\def\labelenumi{(\arabic{enumi})}
\setcounter{enumi}{1}
\item
  Let \(\mathbf{E}\) be a magma with law denoted by \(\top\) . The mapping which associates with \(x \in \mathbf{E}\) the mapping \(\mathbf{A} \mapsto x \top \mathbf{A}\) of the set of subsets of \(\mathbf{E}\) into itself is an action of \(\mathbf{E}\) on \(\mathfrak{P}(\mathbf{E})\) .
\item
  Let E be a set. The identity mapping of \(E^{E}\) is an action of \(E^{E}\) on E, called the canonical action. The corresponding law of action is the mapping \((f, x) \mapsto f(x)\) of \(E^{E} \times E\) into E.
\item
  Let \((\Omega_{i})_{i\in I}\) be a family of sets. For all \(i\in I\) , let \(f_{i}\colon \Omega_{i}\to E^{\mathbb{E}}\) be an action of \(\Omega_{i}\) on E. Let \(\Omega\) be the sum of the \(\Omega_{i}\) (Set Theory, II, § 4, no. 8). The mapping \(f\) of \(\Omega\) onto \(E^{\mathbb{E}}\) , extending the \(f_{i}\) , is an action of \(\Omega\) on E. This allows us to reduce the study of a family of actions to that of a single action.
\item
  Given an action of \(\Omega\) on \(\mathbf{E}\) with law denoted by \(\bot\) , a subset \(\Xi\) of \(\Omega\) and a subset \(\mathbf{X}\) of \(\mathbf{E}\) , \(\Xi \perp \mathbf{X}\) denotes the set of \(\alpha \perp x\) with \(\alpha \in \Xi\) and \(x \in \mathbf{X}\) ; when \(\Xi\) consists of a single element \(\alpha\) , we generally write \(\alpha \perp \mathbf{X}\) instead of \(\{\alpha\} \perp \mathbf{X}\) . The mapping which associates with \(\alpha \in \Omega\) the mapping \(\mathbf{X} \mapsto \alpha \perp \mathbf{X}\) is an action of \(\Omega\) on \(\mathfrak{P}(\mathbf{E})\) , which is said to be derived from the given action by extension to the set of subsets.
\item
  Let \(\alpha \mapsto f_{\alpha}\) be an action of \(\Omega\) on \(\mathbf{E}\) . Let \(g\) be a mapping of \(\Omega'\) into \(\Omega\) . Then the mapping \(\beta \mapsto f_{g(\beta)}\) is an action of \(\Omega'\) on \(\mathbf{E}\) .
\item
  Let \(f: \mathbf{E} \times \mathbf{E} \to \mathbf{E}\) be a law of composition on a set \(\mathbf{E}\) . The mapping \(\gamma: x \mapsto \gamma_x\) (resp. \(\delta: x \mapsto \delta_x\) ) (§ 2, no. 2) which associates with the element \(x \in \mathbf{E}\) left (resp. right) translation by \(x\) is an action of \(\mathbf{E}\) on itself; it is called the left (resp. right) action of \(\mathbf{E}\) on itself derived from the given law. When \(f\) is commutative, these two actions coincide.
\end{enumerate}

The law of left (resp. right) action associated with \(\gamma\) is f (resp. the opposite law to f). The law of right (resp. left) action associated with \(\delta\) is f (resp. the opposite law to f).

Let \(\Omega, \mathbf{E}, \mathbf{F}\) be sets, \(\alpha \mapsto f_{\alpha}\) an action of \(\Omega\) on \(\mathbf{E}\) and \(\alpha \mapsto g_{\alpha}\) an action of \(\Omega\) on \(\mathbf{F}\) . An \(\Omega\) -morphism of \(\mathbf{E}\) into \(\mathbf{F}\) , or mapping of \(\mathbf{E}\) into \(\mathbf{F}\) compatible with the action of \(\Omega\) , is a mapping \(h\) of \(\mathbf{E}\) into \(\mathbf{F}\) such that

\[
g _ {\alpha} (h (x)) = h \left(f _ {\alpha} (x)\right)
\]

for all \(\alpha \in \Omega\) and \(x \in E\) . The composition of two \(\Omega\) -morphisms is an \(\Omega\) -morphism.

Let \(\Omega, \Xi, E, F\) be sets, \(\alpha \mapsto f_{\alpha}\) an action of \(\Omega\) on \(E\) , \(\beta \mapsto g_{\beta}\) an action of \(\Xi\) on \(F\) and \(\phi\) a mapping of \(\Omega\) into \(\Omega'\) . A \(\phi\) -morphism of \(E\) into \(F\) is a mapping \(h\) of \(E\) into \(F\) such that

\[
g _ {\phi (\alpha)} (h (x)) = h (f _ {\alpha} (x))
\]

for all \(\alpha \in \Omega\) and \(x \in \mathbf{E}\) .

